% Régions et zones de disponibilité, selon les conventions des schémas AWS
\documentclass[border=4pt]{standalone}
\usepackage{awsfig}
\newcommand{\groupe}[6]{% x1 y1 x2 y2 couleur/icone titre
  \draw[groupe=#5] (#1,#2) rectangle (#3,#4);
  \node[anchor=north west, inner sep=0pt] at (#1,#2) {#6};
}
\begin{document}
\begin{tikzpicture}[x=1mm,y=1mm]
  % AWS Cloud
  \draw[groupe=Cloud] (0,0) rectangle (250,-90);
  \node[anchor=north west, inner sep=0pt] at (0,0) {\ic[7mm]{g-cloud}};
  \node[groupetitre=Cloud] at (7,0) {AWS Cloud};
  % Région Paris
  \draw[groupe=Region, dash pattern=on 4pt off 2pt] (6,-10) rectangle (178,-84);
  \node[anchor=north west, inner sep=0pt] at (6,-10) {\ic[7mm]{g-region}};
  \node[groupetitre=Region] at (13,-10) {Région Europe (Paris) \code{eu-west-3}};
  % trois AZ
  \foreach \z/\x in {a/12, b/67, c/122}{
    \draw[draw=AZ, line width=.8pt, dash pattern=on 4pt off 2pt] (\x,-20) rectangle (\x+50,-66);
    \node[groupetitre=AZ] at (\x,-20) {Zone de disponibilité \code{eu-west-3\z}};
    \node at (\x+14,-40) {\svc[11mm]{r-dc}{centre de données}};
    \node at (\x+36,-40) {\svc[11mm]{r-dc}{centre de données}};
  }
  \draw[draw=AZ, line width=1.8pt] (62,-58) -- (67,-58);
  \draw[draw=AZ, line width=1.8pt] (117,-58) -- (122,-58);
  \node[remarque, anchor=north west] at (12,-69) {alimentation, refroidissement et réseau propres à chaque zone ;};
  \node[remarque, anchor=north west] at (12,-74) {liaisons privées redondantes entre zones, latence de l'ordre de la milliseconde};
  % incident
  \draw[draw=Security, line width=1.3pt] (158,-39) circle (8.5mm);
  \node[font=\sffamily\small\bfseries, text=Security, anchor=north west, text width=52mm, align=flush left] at (124,-68) {incident dans une zone :\\ les deux autres continuent};
  % autre région
  \draw[groupe=Region, dash pattern=on 4pt off 2pt] (186,-10) rectangle (244,-50);
  \node[anchor=north west, inner sep=0pt] at (186,-10) {\ic[7mm]{g-region}};
  \node[groupetitre=Region, text width=44mm] at (193,-10) {Région Europe\\(Francfort) \code{eu-central-1}};
  \node at (215,-34) {\ic[10mm]{r-dc}\ \ic[10mm]{r-dc}\ \ic[10mm]{r-dc}};
  \node[remarque, text width=56mm, anchor=north west] at (187,-53) {une autre région, totalement indépendante : rien n'y est copié sans votre décision};
  % points de présence
  \node (pop) at (194,-72) {\ic[10mm]{cloudfront}};
  \node[remarque, anchor=west, text width=46mm] at (201,-72) {points de présence : plusieurs centaines de sites, au plus près des utilisateurs};
\end{tikzpicture}
\end{document}
